\section*{Capítulo \ref{ch:g3:continents-and-oceans} --- O globo: continentes e oceanos}

\begin{solution}{exo:g3:continents-and-oceans:1}
A Ásia, a África, a América do Norte, a América do Sul, a Antártida, a
Europa e a Oceania. O maior é a Ásia; o menor é a Oceania.
\end{solution}

\begin{solution}{exo:g3:continents-and-oceans:2}
O Pacífico, o Atlântico, o Índico, o Austral e o Ártico.
O maior é o \omterm{def:g2:water-on-earth:ocean}{oceano} Pacífico; o menor é o \omterm{def:g2:water-on-earth:ocean}{oceano} Ártico.
\end{solution}

\begin{solution}{exo:g3:continents-and-oceans:3}
O \omterm{def:g3:continents-and-oceans:poles}{Polo Norte} fica no alto de um \omterm{def:g3:continents-and-oceans:globe}{globo}; o \omterm{def:g3:continents-and-oceans:poles}{Polo Sul} fica embaixo; o \omterm{def:g3:continents-and-oceans:equator}{equador}
dá a volta inteira pelo meio.
\end{solution}

\begin{solution}{exo:g3:continents-and-oceans:4}
A Europa: norte. A Oceania: sul. A Antártida: sul. A América do Norte:
norte.
\end{solution}

\begin{solution}{exo:g3:continents-and-oceans:5}
O \omterm{def:g2:water-on-earth:ocean}{oceano} Ártico fica em volta do \omterm{def:g3:continents-and-oceans:poles}{Polo Norte}. O \omterm{def:g2:water-on-earth:ocean}{oceano} Austral fica em
volta da Antártida.
\end{solution}

\begin{solution}{exo:g3:continents-and-oceans:6}
Da África à América do Sul: o \omterm{def:g2:water-on-earth:ocean}{oceano} Atlântico. Da Austrália à América do
Sul, indo para o leste: o \omterm{def:g2:water-on-earth:ocean}{oceano} Pacífico. Da África à Austrália: o \omterm{def:g2:water-on-earth:ocean}{oceano}
Índico.
\end{solution}

\begin{solution}{exo:g3:continents-and-oceans:7}
Sam mora perto do \omterm{def:g3:continents-and-oceans:equator}{equador}, onde faz calor o ano inteiro. Bo mora perto de
um \omterm{def:g3:continents-and-oceans:poles}{polo}, onde faz frio o ano inteiro.
\end{solution}

\begin{solution}{exo:g3:continents-and-oceans:8}
$10 - 7 = 3$: o dedo dela cai na terra 3 vezes. É o que podemos esperar:
cerca de 71 quadrados em 100 da Terra são de água; então, mais ou menos 7
vezes em 10, o dedo para no \omterm{def:g2:water-on-earth:ocean}{mar}. (Com só 10 giros, ela também poderia
tirar 6 ou 8: a contagem muda de uma partida para outra.)
\end{solution}

\begin{solution}{exo:g3:continents-and-oceans:9}
Para desenhar o \omterm{def:g3:continents-and-oceans:world-map}{mapa-múndi}, o \omterm{def:g3:continents-and-oceans:globe}{globo} foi cortado ao longo de uma linha que
passa pelo meio do \omterm{def:g2:water-on-earth:ocean}{oceano} Pacífico e aberto até ficar plano. Uma metade do
Pacífico vai parar na beirada da esquerda do \omterm{def:g2:drawing-to-map:map}{mapa}, e a outra metade na
beirada da direita. Na Terra de verdade há um só \omterm{def:g2:water-on-earth:ocean}{oceano} Pacífico: as duas
beiradas do \omterm{def:g2:drawing-to-map:map}{mapa} são a mesma linha no \omterm{def:g3:continents-and-oceans:globe}{globo}.
\end{solution}

\begin{solution}{exo:g3:continents-and-oceans:10}
Tom não tem razão. Do maior para o menor, os \omterm{def:g3:continents-and-oceans:continent}{continentes} são a Ásia, a
África, a América do Norte, a América do Sul, a Antártida, a Europa e a
Oceania: a Antártida é só o quinto, muito menor que a Ásia. O \omterm{def:g3:continents-and-oceans:world-map}{mapa-múndi} o
engana porque um \omterm{def:g2:drawing-to-map:map}{mapa} plano estica muito a terra perto dos \omterm{def:g3:continents-and-oceans:poles}{polos}. A
Antártida, que fica em volta do \omterm{def:g3:continents-and-oceans:poles}{Polo Sul}, é esticada numa faixa comprida
que atravessa toda a parte de baixo do \omterm{def:g2:drawing-to-map:map}{mapa}. No \omterm{def:g3:continents-and-oceans:globe}{globo}, ela é uma massa de
terra redonda, e muito menor que a Ásia.
\end{solution}
